\documentclass[11pt,a4paper]{article}

\usepackage[utf8]{inputenc}
\usepackage[T1]{fontenc}
\usepackage[margin=2.3cm]{geometry}
\usepackage{amsmath}
\usepackage{booktabs}
\usepackage{array}
\usepackage{xcolor}
\usepackage{graphicx}
\usepackage{enumitem}
\usepackage{listings}
\usepackage{hyperref}

\hypersetup{
    colorlinks=true,
    linkcolor=black,
    urlcolor=blue,
}

\lstset{
    basicstyle=\ttfamily\small,
    frame=single,
    breaklines=true,
    columns=fullflexible,
    showstringspaces=false,
}

\title{AIT3002 Midterm Report: QR Code Detection and Decoding}
\author{Vũ Nguyên Đan, 23020351}
\date{\today}

\begin{document}
\maketitle

\section{Method Overview}

The program reads a CSV listing images, detects every QR code in each image,
returns the four-corner quadrilateral of every detection, and optionally
decodes the payload. Everything is written from scratch; no explicit QR library
(\texttt{pyzbar}, \texttt{zxing}, \texttt{cv2.QRCodeDetector}, etc.) and no
deep-learning framework are used.  The only third-party dependencies are
OpenCV (image I/O, thresholding, contour finding, perspective warp) and
NumPy (array math).

The pipeline is:

\begin{enumerate}[leftmargin=1.6em, itemsep=0.2em]
    \item Read the image. If $\max(h,w) > 1024$, resize so the long side is
          $1024$ px (\texttt{INTER\_AREA}). Work in grayscale; precompute
          CLAHE on gray (\texttt{clipLimit}=3.0, \texttt{tileGridSize}=(8,8))
          and a light Gaussian blur of that CLAHE output.
    \item Run the phased binarization cascade in the next section. Each
          phase may append a new candidate pool; pools are merged and passed
          to triplet grouping.
    \item Merge near-duplicate finder candidates across pools, cap the
          merged list, group triplets, estimate corners, deduplicate QRs.
    \item Optionally run a $2\times$ upscale pass on small images when few
          or no QRs are found, and/or a separate dense-grid upscale pass
          when many small finders are present.
    \item Convert each oriented quadrilateral to a $5\%$-expanded
          axis-aligned box for \texttt{output.csv} \texttt{corners}, and
          keep \texttt{oriented\_corners} plus \texttt{finder\_side} for the
          decoder.
    \item If \texttt{--decode=yes}, for every detection run the decoder
          described later in this document.
    \item Write \texttt{output.csv} with one row per detected QR (or one
          empty row for images with zero QRs).
\end{enumerate}

\section{Image preprocessing and detection cascade}

All binarizations use the scaled working image. Phase~1 always runs; later
phases run only when \textbf{no} QR has been assembled yet
(\texttt{len(qrcodes)==0}) and the total candidate count across pools is
$\le 160$ (circuit breaker). Each binary image feeds
\texttt{\_find\_finder\_candidates}; if the pool has $\ge 3$ candidates,
\texttt{\_scan\_line\_verify\_np} must pass on at least two of four directions
when \texttt{strict=True} (Phase~1). Phases~2--5b pass
\texttt{strict=False} (only one direction required) to recover weak
binaries. If verified count falls below three, the unfiltered candidate pool
is kept. Each pool is capped at $90$ entries.

\paragraph{Phase 1 (always).}
\begin{enumerate}[label=(\alph*), leftmargin=1.6em, itemsep=0.15em]
  \item Otsu on blurred CLAHE gray (\texttt{THRESH\_BINARY\_INV}).
  \item Adaptive Gaussian $15\times15$, $C=2$ on the same blur.
  \item Raw gray adaptive Gaussian $11\times11$, $C=2$ only if
        small-finder count $\ge 6$ across pools so far, or total finders
        $<6$, or average finder area $<300$ when at least three finders
        exist.
\end{enumerate}
After Phase~1, \texttt{\_assemble\_qrs} runs on all accumulated pools.

\paragraph{Phases 2--5b (only if still no QRs and not too many candidates).}


In code, these five fallback branches are named Phase~2, Phase~3, Phase~4,
Phase~5a, and Phase~5b.


\begin{enumerate}[label=(\alph*), leftmargin=1.6em, itemsep=0.15em]
  \item Adaptive Gaussian $11\times11$ on CLAHE-equalised gray (no extra blur).
  \item Unsharp mask on CLAHE ($1.8\times$ CLAHE $-$ $0.8\times$ Gaussian blur
        $\sigma=3$), then adaptive Gaussian $15\times15$ on a $3\times3$ blur
        of that sharpened image.
  \item Adaptive Gaussian $25\times25$ on blurred CLAHE, then morphological
        close with a $3\times3$ kernel; pool kept only if it has $\ge 3$
        candidates.
  \item Adaptive mean $21\times21$, $C=4$ on blurred CLAHE; pool kept only if
        $\ge 3$ candidates.
  \item Raw gray adaptive Gaussian $7\times7$, $C=2$; pool kept only if
        $\ge 3$ candidates.
\end{enumerate}

\paragraph{Phase 6 (2$\times$ upscale, independent trigger).}
If the original image long side $\le 720$ px and no QRs are found yet, the
original gray image is resized by $2.0\times$. A second trigger is used for
small images: if the long side $\le 640$ px, and
$1\le|\text{QRs}|\le 8$, and some detection has
\texttt{finder\_side}$<22$ px, the original gray image is resized by
$2.0\times$. Up to three adaptive Gaussian binarizations (blocks $11$, and
if many small finders also $9$ and $13$) run on the upscaled image; new QRs
are merged after mapping corners back with \texttt{scale}/\texttt{up}.

\subsection{Dense-grid rescue (additional pass)}

This pass runs \textbf{in addition} to earlier phases when the long side
$\le 1280$ px, there is at least one accumulated pool, and the count of
finder candidates with contour area $<400$ is $\ge 20$. A second trigger:
long side $\le 720$, $1\le|\text{QRs}|\le 10$, and some QR has
\texttt{finder\_side}$<20$ px. The gray image is upscaled by
\texttt{up = min(2.0, 1400 / max(h,w))}; the pass is skipped if
\texttt{up}$<1.4$. On the upscaled image, adaptive Gaussian (block $11$)
finder search runs with relaxed scan-line verification ($\ge 1$ direction);
pools are capped at $180$ candidates. New QRs are merged if not duplicates.

\subsection{Output geometry}

Ground-truth boxes are axis-aligned. Each oriented quadrilateral is
converted to its axis-aligned bounding rectangle, then expanded by
$5\%$ around the centre (\texttt{\_BBOX\_EXPAND = 1.05}) before integer
\texttt{corners} are written. \texttt{oriented\_corners} and
\texttt{finder\_side} are kept for decoding.

\section{QR Detection}

\subsection{Finder Pattern Candidates}

A QR finder pattern is a $7\times7$ module square with dark/light/dark/light/dark rings
in the ratio $1{:}1{:}3{:}1{:}1$. In a binary image it
shows up as three nested contours.

For each binary image the detector:

\begin{enumerate}[leftmargin=1.6em, itemsep=0.2em]
\item Runs \texttt{cv2.findContours} with \texttt{RETR\_TREE} to obtain the
      full hierarchy.
\item Looks for chains of three nested contours whose areas roughly
      satisfy $\text{outer}/\text{middle}\approx 49/25$ and
      $\text{middle}/\text{inner}\approx 25/9$. Tolerances are loose
      (factor $\le 12$) because perspective distortion and binarization
      errors widen the ratios considerably.
\item Rejects the outer contour if
      $\max(b_w,b_h)/\min(b_w,b_h) > 4.5$ on its bounding rectangle.
\item Verifies each survivor with a scan-line test: four lines through the
      centroid (horizontal, vertical, two diagonals) are traced, the
      binary run lengths are compared to the expected
      $1{:}1{:}3{:}1{:}1$ pattern, and the candidate passes if at least
      two of four directions match.
\end{enumerate}

Within each direction, any five consecutive runs are tested; each run
length $r$ must satisfy $|r - e\cdot u| \le 0.85\,u + 1$ with unit
$u = (\sum r)/7$ and expected pattern $e \in \{1,1,3,1,1\}$. This blocks
most non-finder blobs that still satisfy nested-area ratios.

\subsection{Candidate Merging}

All per-strategy pools are concatenated. Merge radius between candidates
$i$ and $j$ is $0.6 \cdot \max(\sqrt{\text{area}_i}, \sqrt{\text{area}_j})$
(\texttt{merge\_factor}=0.6 in \texttt{\_merge\_candidates}). If Euclidean
distance between centres is below that threshold, they belong to one
cluster; greedy processing in descending \texttt{size} keeps one
representative per cluster. The merged list is capped at $300$ entries
before triplet search.

\subsection{Triplet grouping and assignment}

Triplets are enumerated differently by candidate count $n$:
\begin{itemize}[leftmargin=1.4em, itemsep=0.2em]
  \item $n \le 15$: union-find \texttt{\_spatial\_cluster} splits candidates
        by distance ($k\_factor$ 5.0, then 3.1 on large clusters); all
        $\binom{|C|}{3}$ triples inside each cluster (plus a ``tiny cluster''
        rescue pool) are tested.
  \item $n > 15$: for each finder, only pairs drawn from its $K$ nearest
        neighbours are tried, with $K \in \{4,\ldots,12\}$ chosen from $n$
        (smaller $K$ when $n$ is very large).
\end{itemize}

\texttt{\_check\_triplet} rejects a triple unless all of the following hold:
\begin{itemize}[leftmargin=1.4em, itemsep=0.2em]
  \item Pairwise $\sqrt{\text{area}}$ ratios between the three finders are
        all $<2.3$ (tolerance parameter 1.3).
  \item Sorted vertex angles: largest angle in $(40^\circ, 140^\circ)$,
        smallest angle $\ge 8^\circ$.
  \item From the largest-angle vertex, the two adjacent sides have length
        ratio $\le 3.0$.
  \item Hypotenuse divided by mean adjacent side lies in $(0.85, 2.2)$.
  \item Each adjacent side length is between $1.6\times$ and $34\times$ the
        mean $\sqrt{\text{finder area}}$ (lower bound rejects touching finders;
        upper bound allows high versions with noisy area estimates).
  \item Timing-pattern sampling along both arms from the right-angle vertex
        passes heuristic transition-count checks relative to expected
        module count.
  \item Estimated version from both area- and bbox-based module sizes lies
        in a permissive range around $[1,40]$.
  \item Arms from the right-angle vertex are not too collinear:
        $|\cos\angle(\vec v_0,\vec v_1)| \le 0.80$.
\end{itemize}

Score (lower is better) matches the code: largest angle $\theta$ at the
right-angle vertex, side-arm ratio $r_{\text{side}}$, hypotenuse ratio
$r_{\text{hyp}}$, and normalised size spread $\sigma_{\text{size}}$:
\[
s =
2.5\,|\theta - 90^\circ|
+ 100\,(r_{\text{side}} - 1)
+ 100\,|r_{\text{hyp}} - 1.414|
+ 60\,\sigma_{\text{size}} .
\]
The middle term uses $(r_{\text{side}} - 1)$ without an absolute value,
as in \texttt{main.py}.

Post-filters: when $n \ge 20$, triplets whose legs exceed
$1.8\times$ the 40th percentile of nearest-neighbour distances are dropped
(modal filter against cross-QR triplets). Triplets whose estimated quad
contains more than three finder centres in total are hard-rejected. Finders
already used in accepted triplets can be re-grouped in a second pass when
$3 \le \text{unused} \le 45$.

\subsection{Corner estimation}

From the right-angle vertex and the two opposite finder centres,
\texttt{\_best\_module\_size} picks a module $m$ consistent with
finder-to-finder distances. Outer corners move outward by
\texttt{offset} $= 4m$ along the two arm directions (not $3.5m$: the
extra half-module accounts for label padding). BR blends $0.6\times$
parallelogram completion with $0.4\times$ an extension estimate.

\section{QR Decoding}

Decoding is optional and is controlled by \texttt{--decode=yes} or
\texttt{--decode=no} (default \texttt{yes}). When off, decoding is skipped
and the \texttt{content} field is left empty; when on there is no per-image
time budget. \texttt{QRDecoder.decode} tries two corner orderings (CW mirror
then CCW) and four cyclic rotations of each ($2\times4=8$ attempts) before
giving up on one detection.

\subsection{Perspective Rectification}

The oriented four-corner quadrilateral returned by the detector is
warped to a square of side
$W = \operatorname{clip}(3 \times \text{side}_{\text{px}}, 240, 960)$
with \texttt{cv2.getPerspectiveTransform} and
\texttt{cv2.warpPerspective}. This first warp is intentionally generous
because the detector's quadrilateral is known to overshoot the real QR
edge by somewhere between $0$ and $1.5$ modules.

\subsection{Finder-Pattern Refinement}

Inside the warped image the detector re-runs its finder-pattern search
(using three thresholding variants to cope with brightness loss from
re-sampling). The three finders closest to the canonical corners
$(0,0)$, $(W,0)$, $(0,W)$ are selected; from their centres the module
size $m$ and version $V$ are recovered with
\[
m = \frac{\sqrt{A}}{7}, \qquad
V = \frac{N - 17}{4}, \qquad
N = \operatorname{round}\!\left(\tfrac{d_{\text{mean}}}{m}\right) + 7 .
\]
A second perspective warp maps the three finder centres to their
canonical module coordinates
$(3.5,\,3.5)$, $(N-3.5,\,3.5)$, $(3.5,\,N-3.5)$ at exactly $7$ pixels
per module, so sampling can be done on a fixed integer grid. If the
refinement step fails, the decoder falls back to timing-pattern version
estimation combined with a brute-force search over margin values
(in order) $\{1.0,\,0.5,\,1.5,\,0.0\}$ modules.

\subsection{Bit matrix sampling}

After finder refinement, each module is $7\times7$ pixels; the bit value
uses the centre pixel of that block compared to a threshold map from
\texttt{\_adaptive\_bit\_threshold}: global Otsu on the sample grid (clamped
to $[30,225]$), plus a box blur whose kernel size is
$\max(3,\lfloor N/4\rfloor)$ (odd), blended as $0.7\times\text{local} +
0.3\times\text{global}$. On the pre-refinement fallback path,
\texttt{cv2.remap} samples the first warp on a regular grid; quiet-zone
margin is tried in order $\{1.0, 0.5, 1.5, 0.0\}$ modules.

Before committing to a sampled matrix the decoder checks that at least
two of the three finder-corner $7\times7$ blocks actually look like a
finder (dark outer ring, dark inner $3\times3$). This catches
misaligned samplings that would otherwise decode into nonsense.

\subsection{Format Information}

The 15-bit format string is read from two redundant locations on the
module grid. After XOR with the fixed mask \texttt{0x5412} the 32 valid
code words are compared by Hamming distance; the closest is kept if
its distance is $\le 3$. The decoded bits give the error-correction
level ($\{L, M, Q, H\}$) and mask index ($0{-}7$).

\subsection{Zig-Zag Extraction and Unmasking}

Data is extracted with the standard zig-zag walker: columns are
traversed in pairs right-to-left, alternating bottom-up and top-down;
column 6 (the vertical timing column) is skipped; modules flagged as
function (finders, separators, timing, format, alignment, version
info, dark module) are skipped. For each data module, the selected
mask function is XOR-ed with the bit value.

\subsection{Reed--Solomon Error Correction}

Error correction is done over $GF(2^8)$ with primitive polynomial
\texttt{0x11D}. The implementation comprises:

\begin{itemize}[leftmargin=1.4em, itemsep=0.2em]
  \item Precomputed log / anti-log tables for $GF(256)$ arithmetic.
  \item Syndrome evaluation.
  \item \textbf{Berlekamp--Massey} for the error-locator polynomial.
  \item Chien search for error positions.
  \item \textbf{Forney algorithm} for error magnitudes.
\end{itemize}

For versions with multiple blocks, codewords are first de-interleaved
into their blocks according to the $(c, d, e)$ triples from the ISO
18004 Appendix-I tables, each block is corrected independently, and
the data portions are concatenated.

\subsection{Data Decoding}

The bit stream is parsed by four-bit mode indicators:

\begin{center}
\begin{tabular}{@{}lll@{}}
\toprule
Mode & Indicator & Packing \\
\midrule
Numeric      & \texttt{0001} & 3 digits / 10 bits, 2 / 7, 1 / 4 \\
Alphanumeric & \texttt{0010} & pair / 11 bits, single / 6 \\
Byte         & \texttt{0100} & 8 bits/byte, decoded as UTF-8 with
                                 Latin-1 / Shift-JIS fallback \\
Kanji        & \texttt{1000} & 13 bits/char, Shift-JIS \\
\bottomrule
\end{tabular}
\end{center}

ECI designators, FNC1 markers and structured-append headers are
skipped rather than honoured, which is sufficient for the vast
majority of QRs in the data set.

\section{Libraries and Algorithms}

Third-party dependencies:

\begin{itemize}[leftmargin=1.4em, itemsep=0.15em]
  \item \textbf{OpenCV $\ge$ 4.5} -- \texttt{imread},
        \texttt{cvtColor}, \texttt{createCLAHE},
        \texttt{adaptiveThreshold}, \texttt{threshold}
        (\texttt{OTSU}), \texttt{findContours},
        \texttt{boundingRect}, \texttt{morphologyEx},
        \texttt{getPerspectiveTransform},
        \texttt{warpPerspective}, \texttt{remap},
        \texttt{resize}. No \texttt{QRCodeDetector},
        \texttt{detectAndDecode}, or related QR-specific APIs.
  \item \textbf{NumPy $\ge$ 1.20} -- numerical arrays.
  \item Standard library: \texttt{argparse}, \texttt{csv},
        \texttt{multiprocessing}, \texttt{concurrent.futures}, etc.
\end{itemize}

\noindent
\texttt{requirements.txt} lists only OpenCV and NumPy. The header comment in
\texttt{main.py} mentions SciPy; no SciPy import is used at runtime.

Core algorithms, all implemented from scratch:
contour-hierarchy finder detection,
scan-line $1{:}1{:}3{:}1{:}1$ verification,
triplet geometric scoring,
adaptive $k$-nearest-neighbour triplet enumeration for large $n$,
spatial union--find clustering for small $n$,
dense-grid modal and phantom triplet filters,
dense-QR upscale rescue,
perspective rectification,
adaptive Otsu,
BCH(15, 5) format-code decoding,
GF($2^8$) arithmetic,
Berlekamp--Massey, Chien search, Forney,
zig-zag data extraction, four-mode data parser.

No machine-learning model is trained or used.

\section{Results on the Public Sets}

Runs were executed with 12 worker processes on a consumer CPU. Each
measurement below is wall-clock time reported by the program itself
and includes process-pool startup. Table values are from one run per
setting; wall-clock can vary across runs.

\subsection{Detection Only (\texttt{--decode=no})}

\begin{center}
\begin{tabular}{@{}lrrrrrrr@{}}
\toprule
Split & Images & GT QRs & TP & FP & FN & F1 & Time \\
\midrule
train & 1083 & 1979 & 1184 & 154 & 795 & \textbf{0.7139} &
        15.2 s (14 ms/img) \\
valid &  309 &  508 &  359 &  36 & 149 & \textbf{0.7951} &
         6.2 s (20 ms/img) \\
\bottomrule
\end{tabular}
\end{center}

Precision is $0.885$ (train) and $0.909$ (valid); recall is
$0.598$ / $0.707$. In the measured run, wall-clock time is about
$14$\,ms/image on the train split and about $20$\,ms/image on the valid split under the measured
hardware (valid has fewer images so fixed pool overhead is amortised less).

\subsection{Detection + Decoding (\texttt{--decode=yes})}

F1 is unchanged from above (decoding does not alter the predicted
quadrilaterals). The table adds decoding statistics.

\begin{center}
\begin{tabular}{@{}lrrrr@{}}
\toprule
Split & Detected QRs & Decoded & Decode rate & Time \\
\midrule
train & 1338 & 280 & 20.9\,\% & 39.9 s (37 ms/img) \\
valid &  395 &  89 & 22.5\,\% & 17.3 s (56 ms/img) \\
\bottomrule
\end{tabular}
\end{center}

The provided ground-truth CSVs leave the \texttt{content} column empty,
so public-split content correctness cannot be measured. In
\texttt{evaluate.py}, \texttt{Content\_Accuracy} is computed over total TP,
which follows \texttt{output\_requirement.md} section 5.4; with empty GT
content on public, this value is not informative for decoder correctness.
Inspection of the decoded strings shows that whenever the
decoder succeeds, the output is meaningful (URLs, VCards, \texttt{TEL:},
\texttt{MAILTO:}, and the test strings ``\texttt{Version 1 QR}'' /
``\texttt{Version 2 QR Code Test Image}'' that appear in several
categories).

\subsection{Representative Outputs}

Successful decodes include:
\begin{itemize}[leftmargin=1.4em, itemsep=0.1em]
  \item \texttt{http://www.sneakyparties.com.au}
  \item \texttt{http://qr.westergasfabriek.nl/480}
  \item \texttt{BEGIN:VCARD\textbackslash n VERSION:2.1 \textbackslash n
        N:Gump;Forrest ...}
  \item \texttt{ANA080514DHGE38V4TBT246927890 2070031}
  \item \texttt{TEL:+4992815939305}
\end{itemize}

Typical detection failures:
\begin{itemize}[leftmargin=1.4em, itemsep=0.1em]
  \item Heavy motion blur --- the scan-line test stops validating the
        $1{:}1{:}3{:}1{:}1$ pattern even when all three finders are
        visually present.
  \item Extreme perspective where two finders almost overlap; the
        right-angle test fails.
  \item Partially occluded QRs where one of the three finders is
        physically missing (the pipeline requires three).
  \item Very small QRs in \texttt{lotsimage} (module size below
        $\sim 1.5$ px after downscaling) even after the upscale rescue.
\end{itemize}

Typical decoding failures (given that detection succeeded):
\begin{itemize}[leftmargin=1.4em, itemsep=0.1em]
  \item Finder-pattern refinement picks a spurious finder (text or a
        reflection mimicking one), which throws the sampling grid off by
        a full module.
  \item High noise + curved surfaces: the format code survives but the
        data portion has too many bit errors for Reed--Solomon to
        correct.
  \item Version $\ge 7$: version info block is not used for cross-check,
        so if timing-based version estimation is wrong by one, the
        grid is mis-sampled.
\end{itemize}

\subsection{Speed}

Detection is near the $20$\,ms/image target on the training split.
Decoding is slower because each detection tries eight corner orderings,
then finder refinement; if that fails, up to four versions are tried and
for each version up to four quiet-zone margins are tested before giving
up. Up to twelve (EC level, mask) pairs are tried per matrix.
Decoding has no time cap and can be disabled with \texttt{--decode=no}.

\subsection{Reproducibility}

The pipeline is deterministic: no random sampling and no randomized seed.
\texttt{ProcessPoolExecutor.map} preserves input order, so for fixed code
and input CSV order, output row ordering is deterministic.

\section{Analysis and Discussion}

\subsection{Difficulties}

\begin{itemize}[leftmargin=1.4em, itemsep=0.2em]
  \item Dense grid images (often filenames containing \texttt{lotsimage})
        produce large finder candidate sets. Without spatial clustering
        plus nearest-neighbour triplet pruning, exhaustive triplet search
        would be too slow; without the modal and phantom filters, many
        cross-QR false triplets would survive.
  \item The detector's quadrilateral does not have a fixed margin
        around the real QR edge; it varies between $0$ and
        $\sim 1.5$ modules depending on which rescue path produced it.
        Decoding only became reliable after the second, finder-anchored
        warp was added.
  \item Without ground-truth content strings there is no automatic
        signal for decoder regressions; only the detection F1 and the
        count of non-empty content rows were available during
        development.
\end{itemize}

\subsection{Limitations}

\begin{itemize}[leftmargin=1.4em, itemsep=0.2em]
  \item Requires all three finder patterns: any QR with one finder
        fully occluded is missed.
  \item The format/mask search is a brute-force fallback when BCH
        decoding fails; it can occasionally confirm a wrong mask
        because a few mask+EC combinations happen to yield plausible
        mode indicators.
  \item Version $\ge 7$ version-info codewords are not decoded, so an
        off-by-one version estimate on large QRs goes uncaught.
  \item Kanji mode is implemented but not tested (no Kanji
        QRs in the public set).
\end{itemize}

\subsection{Possible Improvements}

\begin{itemize}[leftmargin=1.4em, itemsep=0.2em]
  \item Two-finder rescue: when two clean finders are found but the
        third is missing, synthesize a candidate from timing-pattern
        extrapolation.
  \item Use detected alignment-pattern centres to refine the
        perspective warp for versions $\ge 2$, reducing the
        misalignment that currently dominates the decoding failure
        cases.
  \item Decode version info for $V \ge 7$ to eliminate the
        timing-based off-by-one.
  \item Add a lightweight quality gate on the decoded text (UTF-8
        validity, printable ratio) before accepting it, to suppress
        the rare but present false positives from brute-force mask
        search.
  \item Port the critical inner loops of the scan-line verifier and
        zig-zag walker to Cython or Numba to drop detection time
        further.
\end{itemize}

\section{Overfitting and Generalization}

The system is a hand-written rule-based pipeline; no model parameters
are learned from the public training set. All thresholds (finder-size
limits, triplet-scoring weights, morphological kernel sizes, adaptive
block sizes, etc.) were chosen from first principles or from the QR
specification. The numbers that depend on the data are:

\begin{enumerate}[leftmargin=1.4em, itemsep=0.15em]
  \item Numeric constants in the cascade (pool caps $90$/$300$/$180$,
        merge factor $0.6$, circuit breaker $160$, dense triggers with
        small-finder area threshold $400$, upscale limits, modal-filter
        percentile $40$ and multiplier $1.8$, etc.).
  \item The $5\%$ axis-aligned box expansion (\texttt{\_BBOX\_EXPAND =
        1.05}) chosen from a sweep on the public training labels.
  \item Scan-line ratio tolerance (default $0.85$ in
        \texttt{\_check\_ratio\_runs}) and triplet scoring weights.
\end{enumerate}

These are hand-set constants, not weights from gradient-based training.
Behaviour on a private test set can differ if image statistics differ,
but nothing in the pipeline memorises per-image labels from the public CSVs.

Re-running the evaluation on the train or valid split yields
byte-identical \texttt{output.csv} files (no randomness, no
timing-dependent behaviour), so any cross-run variation in reported
F1 will come from changes to the program itself, not from
stochasticity.

\section{How to Run}

\begin{lstlisting}
pip install -r requirements.txt
python main.py --data qr/public_train.csv          # default: decode=yes
python main.py --data qr/public_valid.csv --decode=no
python evaluate.py qr/output_valid.csv output.csv
\end{lstlisting}

The program always writes \texttt{output.csv} next to \texttt{main.py}.
The workers count is auto-capped at $\min(12, \text{cpu\_count})$.

\end{document}
